\section{Why neither hypothesis can be dropped}\label{ch17:sec:counterexamples}
As in Section~\ref{ch12:sec:hypotheses}, we remove one hypothesis at a time, keep the other, and look for an example in which the conclusion fails.

\subsection*{Without total boundedness}
Give $\NN$ the discrete metric of Section~\ref{ch17:sec:balls}, in which different numbers are at distance $1$. This space is complete. To see why, take a Cauchy sequence in it and apply the Cauchy condition with the tolerance $1/2$. It supplies a threshold $N$ such that any two terms from $x_N$ onward are at distance less than $1/2$. Distances in this space are only ever $0$ or $1$, so all these distances are $0$, and every term from $x_N$ onward equals $x_N$. Such a sequence converges to $x_N$, a point of the space: for every $\eps>0$ and every $n\ge N$, we have $d(x_n,x_N)=0<\eps$.

The space is not totally bounded, as Section~\ref{ch17:sec:balls} showed. And the conclusion of the theorem fails. Consider the sequence $x_n=n$, that is, $0,1,2,\ldots$. In any subsequence $x_{n_1},x_{n_2},\ldots$ the indices increase, so the terms $n_1,n_2,\ldots$ are different numbers, and any two of them are at distance $1$. For the tolerance $1/2$, no threshold can work, so the subsequence is not a Cauchy sequence. Since every convergent sequence is a Cauchy sequence (Proposition~\ref{ch11:prop:convergent-cauchy}), the subsequence does not converge. In the proof, it is Step~1 that breaks down: no finite collection of balls of radius $1/4$ covers the space, so the selection cannot even begin.

\subsection*{Without completeness}
The interval $(0,1)$, with distance $|x-y|$, is totally bounded. The centers $0$ and $1$ used for $[0,1]$ are not allowed here, because the centers must be points of the space, so we use midpoints instead. Given $r>0$, choose a positive integer $N$ with $1/(2N)<r$, and split $[0,1]$ into the $N$ pieces $\bigl[\tfrac{j-1}N,\tfrac jN\bigr]$ for $j=1,\ldots,N$. Their midpoints $\tfrac{2j-1}{2N}$ lie strictly between $0$ and $1$, so they are points of $(0,1)$. Every $x\in(0,1)$ lies in one of the pieces, and is therefore at distance at most $1/(2N)<r$ from its midpoint.

But $(0,1)$ is not complete (Section~\ref{ch11:sec:complete}), and the conclusion of the theorem fails. Consider the sequence $x_n=1/(n+2)$, that is, $\tfrac12,\tfrac13,\tfrac14,\ldots$, whose terms all lie in $(0,1)$. On the real line it converges to $0$: given $\eps>0$, choose an integer $N>1/\eps$; then every $n\ge N$ gives $x_n=\tfrac1{n+2}<\tfrac1N<\eps$.

Every subsequence $x_{n_1},x_{n_2},\ldots$ also converges to $0$ on the real line. The indices increase by at least one at every step, starting from $n_1\ge0$, so induction gives $n_k\ge k-1$ for every $k\ge1$. Hence, with $\eps$ and $N$ as above, every $k\ge N+1$ has $n_k\ge N$ and therefore $x_{n_k}<\eps$. Now suppose some subsequence converged to a point $L$ of $(0,1)$. Distances in $(0,1)$ are distances on the real line, so the subsequence would converge to $L$ on the real line too. It also converges to $0$ there, so uniqueness of limits (Proposition~\ref{thm:limitunique}) gives $L=0$, which is not a point of $(0,1)$. This contradiction shows that no subsequence converges in $(0,1)$. In the proof, Steps~1--4 go through, and it is Step~5 that fails: the Cauchy subsequence exists, but the point it crowds around is missing from the space.

\subsection*{Two hypotheses, two different jobs}
The two examples fail for different reasons. In the discrete metric on $\NN$, every Cauchy sequence has a limit, but the space cannot be described by finitely many small balls, so there is no way to select terms that crowd together. In $(0,1)$, the selection works perfectly, but the point that the selected terms approach is not in the space. Total boundedness produces a Cauchy subsequence, and completeness gives it a limit. When you meet a new hypothesis, it is worth asking not only ``What does it mean?'' but also ``Which step fails without it?''

\subsection*{Check what you have learned, not only what you can name}
How can you tell whether you have learned a new concept well enough for the proof at hand? Here is a check in four parts, carried out for total boundedness.
\begin{enumerate}
\item \emph{State the problem the concept solves.} Total boundedness provides, for any requested radius, a finite list of balls of that radius that covers the space.
\item \emph{State the definition with its quantifiers in the right order.} The radius comes first. The list, and its length, may depend on the radius, but for that one radius the same list must cover every point.
\item \emph{Test it on examples near the boundary.} The discrete metric on $\NN$ shows that a bound on all distances is not enough to give finite covers by small balls. The interval $(0,1)$ shows that finite covers by small balls do not provide the missing limits.
\item \emph{Return to the exact line of the proof where the concept is used,} here the beginning of each stage of Step~1, and explain it.
\end{enumerate}
If you can recite the definition but cannot explain that line, you have not yet learned the concept well enough for this proof.

The same caution applies to the list of prerequisites itself. An assistant's first list is a proposal, and even in our example it was incomplete. The assistant in Section~\ref{ch17:sec:start} mentioned finite covers, the infinite pigeonhole principle and increasing indices. It did not mention the diameter bound for balls, and it did not explain why nested selections make the chosen terms a Cauchy sequence, yet the proof needs both. Finding such a gap is not a failure of the method; it is the method at work. Add the missing step to the map, learn it, and return to the argument. A list of prerequisites deserves your trust for a particular proof only after you have used it to rebuild that proof, not because an assistant produced it.
